% ================================================
% fig-arch.tex — Архитектура НС ПСПС-ДФ
% ================================================
\usetikzlibrary{fit,arrows.meta,positioning}
\tikzset{
	ndinp/.style={rectangle, rounded corners=3pt, fill=blue!15,
		draw=blue!60, line width=1pt,
		font=\small\bfseries, text centered,
		minimum width=1.8cm, minimum height=0.7cm},
	ndenc/.style={rectangle, rounded corners=3pt, fill=green!15,
		draw=green!50!black, line width=1pt,
		font=\small\bfseries, text centered,
		minimum width=2.0cm, minimum height=0.7cm},
	ndssm/.style={rectangle, rounded corners=3pt, fill=yellow!25,
		draw=orange!70, line width=1pt,
		font=\small\bfseries, text centered,
		minimum width=2.6cm, minimum height=0.8cm},
	ndout/.style={rectangle, rounded corners=3pt, fill=orange!20,
		draw=orange!70!black, line width=1pt,
		font=\small\bfseries, text centered,
		minimum width=1.6cm, minimum height=0.65cm},
	ndphy/.style={rectangle, rounded corners=3pt, fill=red!15,
		draw=red!70, line width=1.5pt,
		font=\small\bfseries, text centered,
		minimum width=2.6cm, minimum height=1.0cm},
	arr/.style={-{Stealth[length=4pt]}, line width=1pt, color=gray!70!black},
	redarr/.style={-{Stealth[length=4pt]}, line width=1.2pt, color=red!70},
	redarrdash/.style={-{Stealth[length=4pt]}, line width=1.2pt,
		color=red!70, dashed},
}

\begin{tikzpicture}[node distance=0.35cm and 0.4cm, scale=0.8]
	
	%% ── ВХОДЫ ────────────────────────────────────────────────
	\node[ndinp] (IMU) {БИНС(9)};
	\node[ndinp, below=of IMU] (GPS) {ГНСС(3)+m(1)};
	\node[ndinp, below=of GPS] (RPM) {RPM(4)};
	
	%% ── ОСНОВНАЯ ЦЕПОЧКА ─────────────────────────────────────
	\node[ndenc, right=0.8cm of GPS] (EMB)
	{\begin{tabular}{c}Linear\\Embedding\\$d=192$\end{tabular}};
	\node[ndenc, right=0.5cm of EMB] (LN1) {LayerNorm};
	\node[ndssm, right=0.5cm of LN1] (SSM)
	{\begin{tabular}{c}Mamba SSM $\times3$\\{\footnotesize state=16, expand=2}\end{tabular}};
	\node[ndenc, right=0.5cm of SSM] (LN2) {LayerNorm};
	
	%% ── ВЫХОДЫ ───────────────────────────────────────────────
	\node[ndout, right=0.8cm of LN2, yshift=+1.0cm] (Hv) {$\hat{v}$(3)};
	\node[ndout, below=0.25cm of Hv]  (Hq) {$\hat{q}$(4)};
	\node[ndout, below=0.25cm of Hq]  (Ht) {$\hat{\tau}$(4)};
	\node[ndout, below=0.25cm of Ht]  (Ha) {$\hat{a}$(3)};
	
	%% ── ОБВОДКА $\hat{\tau}$ и $\hat{a}$ ────────────────────
	%% inner sep увеличен чтобы рамка не касалась узлов вплотную
	\node[draw=red!70, line width=1.2pt, rounded corners=6pt,
	inner sep=4pt, fit=(Ht)(Ha)] (Hgroup) {};
	
	%% ── ФИЗИКА ───────────────────────────────────────────────
	\node[ndphy, below=0.5cm of EMB] (PHY)
	{\begin{tabular}{c}Физ. модель\\$F=k\sum n_i^2$\\$\tau=B[n_i^2]^\top$\end{tabular}};
	
	%% LOSS выровнен по центру между PHY и Hgroup по горизонтали
	\node[ndphy, below=0.9cm of SSM, xshift=0.5cm] (LOSS)
	{$\mathcal{L}_{DF}$};
	
	%% ── СТРЕЛКИ: ИНФЕРЕНС ────────────────────────────────────
	\draw[arr] (IMU.east) -- (EMB.north west);
	\draw[arr] (GPS.east) -- (EMB.west);
	\draw[arr] (RPM.east) -- (EMB.south west);
	\draw[arr] (EMB)  -- (LN1);
	\draw[arr] (LN1)  -- (SSM);
	\draw[arr] (SSM)  -- (LN2);
	\draw[arr] (LN2.east) -- (Hv.west);
	\draw[arr] (LN2.east) -- (Hq.west);
	\draw[arr] (LN2.east) -- (Ht.west);
	\draw[arr] (LN2.east) -- (Ha.west);
	
	%% ── СТРЕЛКИ: ОБУЧЕНИЕ (красные) ─────────────────────────
	\draw[redarr] (RPM.south) |- (PHY.west)
	node[midway, left, font=\footnotesize\bfseries, color=red!70] {обучение};
	
	\draw[redarr] (PHY.east) -- (LOSS.west);
	
	%% Стрелка от обводки (правая сторона Hgroup) к LOSS
	\draw[redarr] (Hgroup.south) |- (LOSS.east);
	
	%% Обратное распространение: от LOSS вверх и к SSM
	\draw[redarrdash] (LOSS.north) -- ++(0,0.5) -| (SSM.south)
	node[near start, right, xshift=-130pt, yshift=2pt,
	font=\footnotesize\bfseries, color=red!70]
	{обратное распространение $\mathcal{L}_{DF}$};
	
\end{tikzpicture}